\clearpage
\SourceLeaf{174}\centerline{\Large CORRESPONDENCE.}

\bigskip
\hrule
\bigskip

\centerline{[LEGENDRE'S THEORY OF PARALLELS.]}

\bigskip
\hrule
\bigskip

\centerline{[\textsc{Gerling} \emph{to} \textsc{Gauss}. \ \emph{Kassel, 11 March 1816.}]}

\smallskip

\{.\ .\ .\ .\ .\ .\ .\ As I close this letter, it occurs to me that I have often forgotten to ask for your judgment on \textsc{Legendre}'s theory of parallels in his élémens de géom. --- He defines the straight line as the shortest path between two points, and uses this definition to prove that the sum of the 3 angles in a triangle can be \textls{no greater} than $2R$. He then proves that it can be \textls{no less} than $2R$ either. This assumes, however, that through a point between the arms of an angle that is $< \tfrac{2}{3}R$ one can always draw a line that intersects both arms. He justifies this assumption in a note (pag. 280, 6th edit. 1806) by letting the lines joining points equidistant from the vertex, one on each arm, catch up with the point and pass it. --- In this last step, however, there seems to me to be the same
error that you showed me I had made during my last stay in Göttingen. He declares it to be \textls{assez évident} and believes that it cannot be made any more rigorous without starting from a different definition of the straight line. --- Afterwards, however, he shows, in yet another way that is new to me and seems rigorous, that the sum of the 3 angles in a triangle must be $= 2R$. It goes roughly as follows. The 2 angles and the included side $A, B, c$ determine the whole triangle, so
\[
C = \varphi(A,B,c).
\]
